\documentclass{stex}

\libusepackage{naproche}
\libinput{preamble}

\begin{document}
\begin{smodule}{primes.ftl}

\importmodule[libraries/arithmetics]{divisibility.ftl}
\importmodule[libraries/arithmetics]{modular-arithmetics.ftl}
\importmodule[libraries/arithmetics]{euclid-division-theorem.ftl}

\symdef{Primes}{\mathbb P}

\symdecl*{prime}
\symdecl*{compound}
\symdecl*{prime number}
\symdecl*{coprime}
\symdecl*{relatively prime}
\symdecl*{mutually prime}

\begin{definition*}[forthel,for={prime,compound,prime number}]
  Let $n$ be a natural number.
  $n$ is \emph{prime} iff $n \Greater \One$ and $n$ has no nontrivial divisors.
  Let $n$ is \emph{compound} stand for $n$ is not prime.
  Let a \emph{prime number} stand for a natural number that is prime.
\end{definition*}

\begin{definition*}[forthel,for=Primes]
  $\emph{\Primes}$ is the class of all prime numbers.
\end{definition*}

\begin{proposition*}[forthel]
  Let $n$ be a natural number such that $n \Greater \One$.
  Then $n$ is prime iff every divisor of $n$ is a trivial divisor of $n$.
\end{proposition*}

\begin{proposition*}[forthel]
  Let $n$ be a natural number such that $n \Greater \One$.
  Then $n$ has a prime divisor.
\end{proposition*}
\begin{proof}[forthel]
  Define $\Phi \DefEq \SClass{n'}{\Naturals}{\text{ if }n' \Greater \One\text{ then }n'\text{ has a prime divisor}}$.

  Let us show that for every $n' \Elem \Naturals$ if $\Phi$ contains all
  predecessors of $n'$ then $\Phi$ contains $n'$.
    Let $n' \Elem \Naturals$.
    Assume that $\Phi$ contains all predecessors of $n'$.
    We have $n' \Eq \Zero$ or $n' \Eq \One$ or $n'$ is prime or $n'$ is composite.

    \begin{case}{$n' \Eq \Zero$ or $n' \Eq \One$.} \end{case}

    \begin{case}{$n'$ is prime.} \end{case}

    \begin{case}{$n'$ is composite.}
      Take a nontrivial divisor $m$ of $n'$.
      Then $\One \Less m \Less n'$.
      $m$ is contained in $\Phi$.
      Hence we can take a prime divisor $p$ of $m$.
      Then we have $p \Divides m \Divides n'$.
      Thus $p \Divides n'$.
      Therefore $p$ is a prime divisor of $n'$.
    \end{case}
  End.

  Thus every natural number belongs to $\Phi$ (by \sr{induction II}{induction}).
\end{proof}

\begin{definition*}[forthel,for={coprime,relatively prime,mutually prime}]
  Let $n, m$ be natural numbers.
  $n$ and $m$ are \emph{coprime} iff for all nonzero natural numbers $k$ such that $k \Divides n$ and $k \Divides m$ we have $k \Eq \One$.
  Let $n$ and $m$ are \emph{relatively prime} stand for $n$ and $m$ are coprime.
  Let $n$ and $m$ are \emph{mutually prime} stand for $n$ and $m$ are coprime.
  Let $n$ is \emph{prime to $m$} stand for $n$ and $m$ are coprime.
\end{definition*}

\begin{proposition*}[forthel]
  Let $n, m$ be natural numbers.
  $n$ and $m$ are coprime iff $n$ and $m$ have no common prime divisor.
\end{proposition*}
\begin{proof}[forthel]
  \begin{case}{$n$ and $m$ are coprime.}
    Let $p$ be a prime number such that $p \Divides n$ and $p \Divides m$.
    Then $p$ is nonzero and $p \NotEq \One$.
    Contradiction.
  \end{case}

  \begin{case}{$n$ and $m$ have no common prime divisor.}
    Assume that $n$ and $m$ are not coprime.
    Let $k$ be a nonzero natural number such that $k \Divides n$ and $k \Divides m$.
    Assume that $k \NotEq \One$.
    Consider a prime divisor $p$ of $k$.
    Then $p \Divides k \Divides n,m$.
    Hence $p \Divides n$ and $p \Divides m$.
    Contradiction.
  \end{case}
\end{proof}

\begin{proposition*}[forthel]
  Let $n, m$ be natural numbers and $p$ be a prime number.
  If $p$ does not divide $n$ then $p$ and $n$ are coprime.
\end{proposition*}
\begin{proof}[forthel]
  Assume $p \NotDivides n$.
  Suppose that $p$ and $n$ are not coprime.
  Take a nonzero natural number $k$ such that $k \Divides p$ and $k \Divides n$.
  Then $k \Eq p$.
  Hence $p \Divides n$.
  Contradiction.
\end{proof}

\begin{proposition*}[forthel]
  Let $n, m$ be natural numbers and $p$ be a prime number.
  If $p \Divides n \Mul m$ then $p \Divides n$ or $p \Divides m$.
\end{proposition*}
\begin{proof}[forthel]
  Assume $p \Divides n \Mul m$.

  \begin{case}{$p \Divides n$.} \end{case}

  \begin{case}{$p \NotDivides n$.}
    Define $\Phi \DefEq \SClass{k}{\Naturals}{k \NotEq \Zero\text{ and }p \Divides k \Mul m}$.
    Then $p \Elem \Phi$ and $n \Elem \Phi$.
    Hence $\Phi$ contains some natural number.
    Thus we can take $a \Elem \Phi$ such that $a \Leq k$ for all $k \Elem \Phi$.

    Let us show that $a$ divides all elements of $\Phi$.
      Let $k \Elem \Phi$.
      Take natural numbers $q, r$ such that $k \Eq (a \Mul q) \Plus r$ and $r \Less a$ (by \sr{Euclid Division Existence}{Euclid's Division Theorem}).
      Indeed $a$ is nonzero.
      Then $k \Mul m
        \Eq ((q \Mul a) \Plus r) \Mul m
        \Eq ((q \Mul a) \Mul m) \Plus (r \Mul m)$.
      We have $p \Divides k \Mul m$.
      Hence $p \Divides ((q \Mul a) \Mul m) \Plus (r \Mul m)$.

      We can show that $p \Divides r \Mul m$.
        We have $p \Divides a \Mul m$.
        Hence $p \Divides (q \Mul a) \Mul m$.
        Indeed $((q \Mul a) \Mul m) \Eq q \Mul (a \Mul m)$. %!
        Take $A \Eq (q \Mul a) \Mul m$ and $B \Eq r \Mul m$. %!
        Then $p \Divides A \Plus B$ and $p \Divides A$.
        Thus $p \Divides B$ (by \sn{divisibility of summands}).
        Indeed $p, A$ and $B$ are natural numbers.
        Consequently $p \Divides r \Mul m$.
      End.

      Therefore $r \Eq \Zero$.
      Indeed if $r \NotEq \Zero$ then $r$ is an element of $\Phi$ that is less than $a$.
      Hence $k \Eq q \Mul a$.
      Thus $a$ divides $k$.
    End.

    Then we have $a \Divides p$ and $a \Divides n$.
    Hence $a \Eq p$ or $a \Eq \One$.
    Thus $a \Eq \One$.
    Indeed if $a \Eq p$ then $p \Divides n$.
    Then $\One \Elem \Phi$.
    Therefore $p \Divides \One \Mul m \Eq m$.
  \end{case}
\end{proof}

\end{smodule}
\end{document}
